\begin{afterword}
\summaryhead

The post answers a defense of the lottery: that a ticket is a rational purchase of fantasy, a dollar for a day of imagining oneself a millionaire. The author replies that the fantasy is of wealth with a probability near zero, which the buyer can do nothing to bring about, and which arrives without effort.

The lottery is therefore a sink of emotional energy as well as of money. Without it, people might dream about things they could do, such as going to technical school, starting a business or earning a promotion, and in repeating the dream might find a way to do it. The author proposes calling ticket-buying stupid. People buy tickets, the post says, because the brain cannot scale a feeling down by a factor of 0.00000001, and because they treat the expected-utility calculation as one argument among others instead of letting it decide. It ends by listing five steps for overcoming a bias and saying that many people stall at the third, deciding that the bias is bad.

\responsehead

The post's figures are sound and its account of why tickets sell has support. The one-in-100-million factor is within a factor of two of the 2007 jackpot odds for the two large American lotteries. That people cannot give a very small probability a correspondingly small weight is what Kahneman and Tversky had said in 1979, and they used it to explain gambling. The post states it without a source, but that is not where the post goes wrong.

It goes wrong in the argument against the defense. The defenders say the anticipation is worth the dollar. The post's case against them has two parts. The first is that a numerical calculation of expected utility ``ought to override'' the buyer's instincts instead of being weighed against pleasant anticipation. But no calculation is shown, and any expected-utility calculation must put a value on the anticipation itself, which is the point in dispute. The second is that lottery dreams take the place of better ones: without the lottery, people ``maybe'' would dream about school or a promotion, and their dreams ``might'' show them the way. This is offered with no evidence. The research closest to it, on positive fantasies about a desired future, found that such fantasies predicted less effort and worse results, including among graduates looking for jobs. On that evidence, the dream the post would put in the lottery's place is not clearly better.

The closing list of steps presents the dispute as settled. The defenders hold that no bias is involved. Right after asking why so many defend the lottery, the post says that many people ``bog down in step 3,'' ``deciding that the bias is bad,'' which ``by rights should be the easiest.'' That describes the disagreement as a failure to take an easy step, when whether there is a bias to reject is the question in dispute.

In short: a sound account of why tickets sell, and a case against buying them that depends on an uncomputed calculation and on a hedged claim about better dreams that the research nearest to it does not support.
\end{afterword}
